\section{The theta source and its heat transform}\label{sec:theta}

The question motivating this paper is whether finite Euler operations can
be written explicitly at positive heat time and then reorganized into
expressions whose signs or sizes can be controlled. There are two distinct
issues. A finite algebraic identity must first survive the heat evolution
with all mixed terms and boundaries retained. An estimate must then control
the particular signed expression of interest. The estimates proved here
address positive arithmetic portions of the theta transform. The complete
transform contains a cancellation between those portions and its first
source mode, and this cancellation remains essential at a zero.

The results have three different scopes. Sections~\ref{sec:packets}--%
\ref{sec:controls} give exact identities, positive auxiliary estimates,
and scoped collision controls. Sections~\ref{sec:shifted}--\ref{sec:effective}
give finite approximations with explicit derivative errors, including
effective unbounded-height tails. Section~\ref{sec:finite} consumes those
errors in finite certificates and a propagation theorem, while
Section~\ref{sec:outer} supplies a separate eventual real exclusion
from the published Riemann--Siegel approximation.
The disconnected coverage and outstanding complex-barrier hypothesis
are stated explicitly in Section~\ref{sec:global}. A bound on a positive
packet or on an omitted tail is not a lower bound for the complete
oscillatory field.

We develop the definitions and analytic justifications before introducing
the arithmetic decomposition. The variable $u$ denotes a source coordinate,
$t$ is heat time, $z$ is a complex spectral coordinate, and $x$ denotes a
real spectral coordinate. All source sums start at $n=1$ unless an index set
is stated. No prior investigation or internal notation is used.

\subsection{Theta inversion from Gaussian Fourier analysis}

For $y>0$, define
\begin{equation}\label{eq:theta}
 \Theta(y)=\sum_{n\in\Z}e^{-\pi n^2y}.
\end{equation}
The series, and its derivatives of any fixed order, converge uniformly
on compact subsets of the right half-plane $\Re y>0$. Indeed, on a compact
set with $\Re y\ge c>0$, each differentiated term is bounded by a constant
times $n^{2k}e^{-\pi c n^2}$, a summable sequence.

We use the Fourier convention
$\widehat f(\xi)=\int_\R f(v)e^{-2\pi i v\xi}\dd v$.
For $f(v)=e^{-\pi yv^2}$, differentiation under the integral and integration
by parts show that
$\widehat f'(\xi)=-(2\pi\xi/y)\widehat f(\xi)$.
The Gaussian integral gives $\widehat f(0)=y^{-1/2}$, so
\[
 \widehat f(\xi)=y^{-1/2}e^{-\pi\xi^2/y}.
\]
The periodization $\sum_{n\in\Z}f(v+n)$ has Fourier coefficients
$\widehat f(k)$. Both this series and its Fourier series converge
absolutely and uniformly. Evaluating at $v=0$ therefore proves
\begin{equation}\label{eq:inversion}
 \Theta(y)=y^{-1/2}\Theta(1/y).
\end{equation}
This is the zero-elliptic-variable form of the classical theta
transformation; see \cite{DLMF}. The preceding argument fixes its
normalization without an appeal to conventions for theta functions.

\subsection{The primitive and the physical source}

Set $\lambda_n=\pi n^2$ and define, for real $u$,
\begin{align}
 p_n(u)&=e^u e^{-\lambda_n e^{4u}},&
 \psi(u)&=\sum_{n\ge1}p_n(u),\label{eq:primitive}\\
 \phi_n(u)&=e^u\bigl(2\lambda_n^2e^{8u}-3\lambda_ne^{4u}\bigr)
                 e^{-\lambda_ne^{4u}},&
 \Ph(u)&=\sum_{n\ge1}\phi_n(u).\label{eq:source}
\end{align}
These sums converge locally with every derivative on the whole real line.
For example, differentiating a term introduces a polynomial in
$\lambda_ne^{4u}$; on a compact $u$ interval the remaining exponential is
bounded by $e^{-c n^2}$ for some $c>0$.

Writing $a=\lambda_ne^{4u}$, direct differentiation gives
\[
 p_n'=e^{u-a}(1-4a),\qquad
 p_n''=e^{u-a}(1-24a+16a^2).
\]
Consequently
\begin{equation}\label{eq:source-operator}
 8\phi_n=p_n''-p_n,\qquad
 8\Ph=\psi''-\psi.
\end{equation}
The completed primitive
\begin{equation}\label{eq:P}
 P(u)=e^u\Theta(e^{4u})=e^u+2\psi(u)
\end{equation}
is even by \eqref{eq:inversion}: substituting $y=e^{4u}$ gives
$P(u)=e^{-u}\Theta(e^{-4u})=P(-u)$. Thus
\begin{equation}\label{eq:reflection}
 \Ph=\frac{P''-P}{16},\qquad
 \Ph(u)=\Ph(-u),\qquad \psi'(0)=-\frac12.
\end{equation}
The last equality follows from $0=P'(0)=1+2\psi'(0)$.
The $e^u$ term in $P$ is annihilated by $\partial_u^2-1$; it must nevertheless
be retained when determining this endpoint constant.

For $u\ge0$, $a\ge\pi>3/2$, so
$\phi_n(u)=e^{u-a}a(2a-3)>0$. Reflection then gives
$\Ph(u)>0$ on the whole real line. Positivity of the source does not
by itself imply positivity of its oscillatory transform.

\begin{lemma}[Decay and differentiation]\label{lem:regularity}
For each real $t$, $w_t(u)=e^{tu^2}\Ph(u)$ is an even Schwartz function on
$\R$. On the positive half-line, the same rapid decay, with every derivative,
holds for $e^{tu^2}\sum c_n\phi_n(u)$ whenever $(c_n)$ is bounded. All of
these bounds are uniform when $t$ ranges over a compact interval. The
unweighted source sums are real analytic on $\R$.
\end{lemma}
\begin{proof}
Each fixed derivative of a source term is a finite sum of terms
$C e^{bu}n^{2k}e^{-\pi n^2e^{4u}}$. For $u\ge0$,
\[
 \sum_{n\ge1}n^{2k}e^{-\pi n^2e^{4u}}
 \le e^{-\pi e^{4u}}
       \sum_{n\ge1}n^{2k}e^{-\pi(n^2-1)}.
\]
The constant series is finite. Differentiation of $e^{tu^2}$ introduces
only polynomial factors in $u$ and $t$. Hence every weighted source
derivative is bounded on $u\ge0$ by a finite sum of expressions
$C(1+u)^j\exp(Tu^2+bu-\pi e^{4u})$, uniformly for $|t|\le T$.
Such expressions decay faster than every power of $u$, and remain
integrable after multiplication by $e^{Yu}$ for any fixed $Y$.
For the complete source, evenness extends these bounds to $u<0$.
For real analyticity, extend $u$ to a sufficiently small complex
neighborhood of any fixed real point. There $\Re(e^{4u})$ has a positive
lower bound, and the source series converges normally. It therefore defines
a holomorphic function on that neighborhood.
\end{proof}

\subsection{The heat family and its normalization}

Define
\begin{align}
 F(t,z)&=\int_0^\infty e^{tu^2+izu}\Ph(u)\dd u,\label{eq:F}\\
 H(t,z)&=\int_0^\infty e^{tu^2}\Ph(u)\cos(zu)\dd u
       =\frac{F(t,z)+F(t,-z)}2.\label{eq:H}
\end{align}
Lemma~\ref{lem:regularity} justifies differentiation under these integrals
on every compact complex $(t,z)$ set. Thus they are entire in both variables,
and
\begin{equation}\label{eq:heat}
 H_t=-H_{zz},\qquad F_t=-F_{zz}.
\end{equation}
For real $t$, $H(t,z)$ is even in $z$ and satisfies
$H(t,\overline z)=\overline{H(t,z)}$. In particular it is real on the real
axis. At complex $z$, the symmetrization in \eqref{eq:H} cannot be replaced
by a real part. Since $w_t$ is even,
\begin{equation}\label{eq:wholeline}
 H(t,z)=\frac12\int_\R w_t(u)e^{izu}\dd u.
\end{equation}
For real $x$, $H(t,x)$ and all its $x$ derivatives are rapidly decreasing,
uniformly on compact time intervals. This also follows directly by repeated
integration by parts in \eqref{eq:wholeline}.

To identify the time-zero normalization, start with
$\zeta(s)=\sum_{n\ge1}n^{-s}$ for $\Re s>1$. Substitution in the defining
gamma integral gives
\[
 \pi^{-s/2}\Gamma(s/2)\zeta(s)
 =\frac12\int_0^\infty (\Theta(y)-1)y^{s/2-1}\dd y.
\]
Split the integral at $1$, apply \eqref{eq:inversion} below $1$, and set
$y=1/v$ in the remaining theta term. The elementary part is
$1/(s-1)-1/s=1/[s(s-1)]$. Therefore
\begin{equation}\label{eq:mellin-split}
 \pi^{-s/2}\Gamma(s/2)\zeta(s)
 =\frac1{s(s-1)}+\frac12\int_1^\infty(\Theta(y)-1)
 \bigl(y^{s/2-1}+y^{(1-s)/2-1}\bigr)\dd y.
\end{equation}
The integral on the right is entire in $s$. Multiplying by $s(s-1)/2$
continues the completed function
\[
 \xi(s)=\frac{s(s-1)}2\pi^{-s/2}\Gamma(s/2)\zeta(s)
\]
to an entire function. With $y=e^{4u}$, the integral in
\eqref{eq:mellin-split} becomes
$16\int_0^\infty\psi(u)\cosh((2s-1)u)\dd u$. At
$s=(1+iz)/2$ this gives
\[
 \xi\bigl((1+iz)/2\bigr)
 =\frac12-(1+z^2)\int_0^\infty\psi(u)\cos(zu)\dd u.
\]
On the other hand, integrating \eqref{eq:source-operator} twice and using
$\psi'(0)=-1/2$ gives exactly the same expression for $8H(0,z)$.
We have proved
\begin{equation}\label{eq:xi}
 H(0,z)=\frac18\xi\bigl((1+iz)/2\bigr).
\end{equation}
The Riemann hypothesis is consequently the assertion that all zeros of
$H(0,\cdot)$ are real. Nothing below assumes this assertion.
